% Complete reviewed Unit 12; editorial master: chapters/ch12/answer-key.md
% Markdown SHA-256: 42493d3250a9403ceb68138c0acee3150b76a70deaa5532b65a01e6bb8c81ae6
% One-use initial assembly followed by complete content and page review.
\newcommand{\runningtitle}{Unit 12. Answer key}
% Layout inherited from the reviewed previous course; set a running title per unit.
% This is a small hand-maintained TeX layout, not a conversion pipeline.
\RequirePackage{fix-cm}
\documentclass[11pt,twoside]{article}
\usepackage{fontspec}
\setmainfont{Cambria}[
  Renderer=HarfBuzz,
  BoldFont=cambriab.ttf,
  ItalicFont=cambriai.ttf,
  BoldItalicFont=cambriaz.ttf
]
\setmonofont{lmmono10-regular.otf}[BoldFont=lmmonolt10-bold.otf]
\usepackage[paperwidth=165mm,paperheight=241mm,inner=21mm,outer=19mm,
  top=22mm,bottom=20mm,headheight=14pt,headsep=7mm,footskip=10mm]{geometry}
\usepackage{xcolor,array,booktabs,longtable,tabularx}
\usepackage{enumitem,fancyhdr}
\usepackage[unicode,hidelinks]{hyperref}
\definecolor{ink}{HTML}{000000}
\definecolor{muted}{HTML}{4C5558}
\setlength{\parindent}{0pt}
\setlength{\parskip}{3.5pt}
\AtBeginDocument{\fontsize{11}{14.3}\selectfont}
\setlength{\tabcolsep}{4pt}
\renewcommand{\arraystretch}{1.10}
\setlist[itemize]{leftmargin=1.3em,label=\textbullet,itemsep=2pt,topsep=3pt,parsep=0pt}
\setlist[enumerate]{leftmargin=1.8em,itemsep=3pt,topsep=4pt,parsep=0pt}
\widowpenalty=10000
\clubpenalty=10000
\emergencystretch=1em
\raggedbottom
\pagestyle{fancy}
\fancyhf{}
\providecommand{\runningtitle}{Old English}
\fancyhead[LE,RO]{\small\itshape \runningtitle}
\fancyhead[LO,RE]{\small Old English}
\fancyfoot[LE,RO]{\small\thepage}
\renewcommand{\headrulewidth}{0.3pt}
\newcommand{\booktitle}[2]{
  {\fontsize{9}{11}\selectfont\color{muted}#1\par}
  \vspace{3pt}
  {\fontsize{20}{24}\selectfont\color{ink}#2\par}
  \vspace{6pt}
}
\newcommand{\source}[1]{\par{\footnotesize\color{muted}#1\par}}
\newcommand{\readingtext}[1]{
  \par\addvspace{5pt}
  \begingroup\leftskip=1.5em
    \fontsize{12}{16}\selectfont #1\par
  \endgroup\addvspace{5pt}
}
\makeatletter
\renewcommand{\section}{\@startsection{section}{1}{\z@}
  {-2.4ex \@plus -0.6ex \@minus -0.2ex}{0.8ex \@plus 0.2ex}
  {\normalfont\fontsize{12}{15}\selectfont\bfseries\color{ink}}}
\renewcommand{\subsection}{\@startsection{subsection}{2}{\z@}
  {-1.8ex \@plus -0.5ex \@minus -0.2ex}{0.5ex \@plus 0.2ex}
  {\normalfont\fontsize{11}{14}\selectfont\bfseries\color{ink}}}
\makeatother

\hyphenpenalty=10000
\exhyphenpenalty=10000
\emergencystretch=3em
\newcommand{\itempara}[2]{\par\addvspace{3pt}\noindent\hangindent=1.8em\hangafter=1\makebox[1.8em][l]{#1.}#2\par}
\newcommand{\reserve}[1]{\par\penalty-100\begingroup\dimen0=\pagegoal\advance\dimen0 by-\pagetotal\ifdim\dimen0<#1\newpage\fi\endgroup}
\begin{document}
\addfontfeatures{Ligatures=TeXOff}
\raggedright
\booktitle{UNIT 12}{Answer key}

% mdparagraph:P001:begin
The key covers all 76 tasks. Cook's three complete sets retain their original prompts and numbering; the answers, translations and explanations are new teaching material. The twelve translated teaching examples are not additional exercises. Equivalent English wording, defensible Old English word order and documented inflectional alternatives are acceptable. A grammatical answer can nevertheless miss a stated target, such as a particular auxiliary, mood or selected paradigm. [1]
% mdparagraph:P001:end

% mdparagraph:P002:begin
In form tasks, supply every labeled position, including positions with identical spellings. In translation and analysis, distinguish morphological tense from the time relation established by context. An English perfect, pluperfect or future does not mechanically require a corresponding Old English auxiliary. In reports, identify the content's time relative to the reporting situation; a subjunctive does not by itself make the report false. [2]
% mdparagraph:P002:end

% mdparagraph:P003:begin
The productive \textbf{habban} construction below has an active subject, an accusative object and a selected past participle. Its finite auxiliary agrees with the subject. The models use bare participles; source usage also permits participles inflected in agreement with the object. A correctly evidenced agreeing alternative is not wrong merely because it differs from the model, but participial agreement is not agreement in person with the finite subject. The twelve indicative construction positions do not constitute twelve new inflected words or exhaust the construction's possible moods. [3]
% mdparagraph:P003:end

% mdparagraph:P004:begin
The \textbf{wesan + gewiten} and \textbf{wesan + sittende} tasks require recognition, not independent composition. Local Cook and reading helps remain available during the tasks they support. Supplied expressions are not tests of unaided recall. An agentless passive does not name the agent; do not restore one merely from a nearby model.
% mdparagraph:P004:end

\clearpage
% mdparagraph:P005:begin
\section*{F. Complete forms and constructions}
% mdparagraph:P005:end

\reserve{9\baselineskip}
\par\begin{minipage}{\linewidth}\raggedright
% mdparagraph:P006:begin
\textbf{F01. Fōt.}
% mdparagraph:P006:end

% item: NEW-U12-F01
% key: KEY-NEW-U12-F01
% mdtable:T01:begin
\par\begingroup
\fontsize{10.5}{13.2}\selectfont
\setlength{\LTleft}{0pt}\setlength{\LTright}{\fill}
\noindent
\begin{tabular}{@{}>{\raggedright\arraybackslash}p{39.46mm}>{\raggedright\arraybackslash}p{39.46mm}>{\raggedright\arraybackslash}p{39.46mm}@{}}
\toprule
\textbf{Case} & \textbf{Singular} & \textbf{Plural} \\
\midrule
Nominative & fōt & fēt \\
Genitive & fōtes & fōta \\
Dative & fēt & fōtum \\
Accusative & fōt & fēt \\
\bottomrule
\end{tabular}
\par\endgroup

% mdtable:T01:end
% mdparagraph:P007:begin
\textbf{Explanation:} The selected root-noun system has \textbf{fēt} in the dative singular and nominative/\allowbreak{}accusative plural. Genitive singular and genitive/\allowbreak{}dative plural retain \textbf{fōt-}. Thus neither vowel can be extended through the whole paradigm. The documented dative singular \textbf{fōte} and occasional nominative/\allowbreak{}accusative plural \textbf{fōtas} are valid alternatives in those cells; the displayed forms give the requested selected table. A genitive plural made from \textbf{fēt-} would miss the assigned alternation. [4]
% mdparagraph:P007:end

\end{minipage}\par\addvspace{4pt}

\par\begin{minipage}{\linewidth}\raggedright
% mdparagraph:P008:begin
\textbf{F02. Brycg and hęre.}
% mdparagraph:P008:end

% item: NEW-U12-F02
% key: KEY-NEW-U12-F02
% mdtable:T02:begin
\par\begingroup
\fontsize{10.5}{13.2}\selectfont
\setlength{\LTleft}{0pt}\setlength{\LTright}{\fill}
\noindent
\begin{tabular}{@{}>{\raggedright\arraybackslash}p{22.13mm}>{\raggedright\arraybackslash}p{23.68mm}>{\raggedright\arraybackslash}p{22.13mm}>{\raggedright\arraybackslash}p{22.68mm}>{\raggedright\arraybackslash}p{22.13mm}@{}}
\toprule
\textbf{Case} & \textbf{Brycg singular} & \textbf{Brycg plural} & \textbf{Hęre singular} & \textbf{Hęre plural} \\
\midrule
Nominative & brycg & brycga & hęre & hęras \\
Genitive & brycge & brycga & hęres & hęra \\
Dative & brycge & brycgum & hęre & hęrum \\
Accusative & brycge & brycga & hęre & hęras \\
\bottomrule
\end{tabular}
\par\endgroup

% mdtable:T02:end
% mdparagraph:P009:begin
\textbf{Explanation:} Feminine \textbf{brycg} retains \textbf{brycg-} before its endings, including accusative singular \textbf{-e}. Masculine \textbf{hęre} has the selected singular in \textbf{-e}, genitive \textbf{hęres}, and plural \textbf{hęr-} before \textbf{-as, -a, -um}. Do not add these endings to the whole nominative as though its final \textbf{e} were retained everywhere. Cook's expanded \textbf{hęrg-/\allowbreak{}hęrig-/\allowbreak{}hęrige-} alternatives are acceptable in the cells to which the learner's table assigns them; they do not authorize arbitrary interchange of stems. A masculine \textbf{-as} plural for \textbf{brycg} would import the wrong system. [4]
% mdparagraph:P009:end

\end{minipage}\par\addvspace{4pt}

\par\begin{minipage}{\linewidth}\raggedright
% mdparagraph:P010:begin
\textbf{F03. Flēon and grētan: complete indicative systems.}
% mdparagraph:P010:end

% item: NEW-U12-F03
% key: KEY-NEW-U12-F03
% mdparagraph:P011:begin
\textbf{Flēon — strong verb, present and preterite indicative.}
% mdparagraph:P011:end

% mdtable:T03:begin
\par\begingroup
\fontsize{10.5}{13.2}\selectfont
\setlength{\LTleft}{0pt}\setlength{\LTright}{\fill}
\noindent
\begin{tabular}{@{}>{\raggedright\arraybackslash}p{44.51mm}>{\raggedright\arraybackslash}p{36.93mm}>{\raggedright\arraybackslash}p{36.93mm}@{}}
\toprule
\textbf{Person and number} & \textbf{Present} & \textbf{Preterite} \\
\midrule
First singular & flēo & flēah \\
Second singular & flīehst & fluge \\
Third singular & flīehð & flēah \\
First plural & flēoð & flugon \\
Second plural & flēoð & flugon \\
Third plural & flēoð & flugon \\
\bottomrule
\end{tabular}
\par\endgroup

% mdtable:T03:end
% mdparagraph:P012:begin
\textbf{Grētan — weak verb, present and preterite indicative.}
% mdparagraph:P012:end

% mdtable:T04:begin
\par\begingroup
\fontsize{10.5}{13.2}\selectfont
\setlength{\LTleft}{0pt}\setlength{\LTright}{\fill}
\noindent
\begin{tabular}{@{}>{\raggedright\arraybackslash}p{44.51mm}>{\raggedright\arraybackslash}p{36.93mm}>{\raggedright\arraybackslash}p{36.93mm}@{}}
\toprule
\textbf{Person and number} & \textbf{Present} & \textbf{Preterite} \\
\midrule
First singular & grēte & grētte \\
Second singular & grētst & grēttest \\
Third singular & grētt & grētte \\
First plural & grētað & grētton \\
Second plural & grētað & grētton \\
Third plural & grētað & grētton \\
\bottomrule
\end{tabular}
\par\endgroup

% mdtable:T04:end
% mdparagraph:P013:begin
\textbf{Explanation:} Contracted \textbf{flēon} has the selected present contrast \textbf{flēo-/\allowbreak{}flīeh-} and strong past \textbf{flēah} beside \textbf{flug-}. Its second-singular preterite uses the plural stem plus \textbf{-e}: \textbf{fluge}, with short \textbf{u}. Weak \textbf{grētan} keeps its long root vowel and forms its past with the dental sequence \textbf{-tt-}. Its third-singular present \textbf{grētt} and first/\allowbreak{}third-singular past \textbf{grētte} must not be confused. Keep all three plural person labels even when their forms coincide. Documented variants remain valid in their proper cells; a strong past invented for \textbf{grētan} does not complete its selected weak system. [5]
% mdparagraph:P013:end

\end{minipage}\par\addvspace{4pt}

\par\begin{minipage}{\linewidth}\raggedright
% mdparagraph:P014:begin
\textbf{F04. Diagnostic indicative forms.}
% mdparagraph:P014:end

% item: NEW-U12-F04
% key: KEY-NEW-U12-F04
% mdtable:T05:begin
\par\begingroup
\fontsize{10.5}{13.2}\selectfont
\setlength{\LTleft}{0pt}\setlength{\LTright}{\fill}
\noindent
% Layout repair: preserve the full geopenian label within its column.
\begin{tabular}{@{}>{\raggedright\arraybackslash}p{18mm}>{\raggedright\arraybackslash}p{23mm}>{\raggedright\arraybackslash}p{22mm}>{\raggedright\arraybackslash}p{26mm}>{\raggedright\arraybackslash}p{23.76mm}@{}}
\toprule
\textbf{Verb} & \textbf{Present second singular} & \textbf{Present third singular} & \textbf{Preterite second singular} & \textbf{Preterite third plural} \\
\midrule
gemētan & gemētst & gemētt & gemēttest & gemētton \\
ādrīfan & ādrīfst & ādrīfð & ādrife & ādrifon \\
geopenian & geopenast & geopenað & geopenodest & geopenodon \\
cȳðan & cȳst & cȳð & cȳðdest & cȳðdon \\
\bottomrule
\end{tabular}
\par\endgroup

% mdtable:T05:end
% mdparagraph:P015:begin
\textbf{Explanation:} Weak \textbf{gemētan} has present \textbf{gemētt} and preterite stem \textbf{gemētt-}, with \textbf{-e} in first and third singular, \textbf{-est} in second singular, and \textbf{-on} in all plural persons. Strong \textbf{ādrīfan} has long \textbf{ī} in the present but short \textbf{i} in the second-singular and plural past stem. Weak-II \textbf{geopenian} uses present \textbf{geopena-} in these singular cells and preterite \textbf{geopenod-}. In selected \textbf{cȳst/\allowbreak{}cȳð}, the dental consonants do not remain as a mechanically concatenated stem and ending; the past retains \textbf{cȳðd-}. \textbf{Cȳðst} follows the documented analogical retained-dental pattern and is a valid alternative, although \textbf{cȳst} matches the selected table. Do not transfer \textbf{ādrāf} into its second-singular past or give \textbf{geopenian} a weak-I past. [5]
% mdparagraph:P015:end

\end{minipage}\par\addvspace{4pt}

\par\begin{minipage}{\linewidth}\raggedright
% mdparagraph:P016:begin
\textbf{F05. Complete command pairs.}
% mdparagraph:P016:end

% item: NEW-U12-F05
% key: KEY-NEW-U12-F05
% mdtable:T06:begin
\par\begingroup
\fontsize{10.5}{13.2}\selectfont
\setlength{\LTleft}{0pt}\setlength{\LTright}{\fill}
\noindent
\begin{tabular}{@{}>{\raggedright\arraybackslash}p{36.04mm}>{\raggedright\arraybackslash}p{41.8mm}>{\raggedright\arraybackslash}p{40.54mm}@{}}
\toprule
\textbf{Verb} & \textbf{Singular command} & \textbf{Plural command} \\
\midrule
flēon & flēoh & flēoð \\
ādrīfan & ādrīf & ādrīfað \\
gemētan & gemēt & gemētað \\
geopenian & geopena & geopeniað \\
\bottomrule
\end{tabular}
\par\endgroup

% mdtable:T06:end
% mdparagraph:P017:begin
\textbf{Explanation:} Each first form addresses one person, each second form several. Preserve the long vowel in \textbf{flēoh}. The weak-II plural \textbf{geopeniað} retains \textbf{i}; \textbf{geopenað} is the selected third-singular indicative, not this plural command. A command may share a spelling with another verbal form, but its function and addressee remain distinct. Substituting indicative \textbf{flīehð} for singular \textbf{flēoh} misses the command target. [5]
% mdparagraph:P017:end

\end{minipage}\par\addvspace{4pt}

\reserve{8\baselineskip}
% mdparagraph:P018:begin
\textbf{F06. Habban with object and past participle.}
% mdparagraph:P018:end

% item: NEW-U12-F06
% key: KEY-NEW-U12-F06
\reserve{11\baselineskip}
% mdparagraph:P019:begin
\textbf{Present indicative construction.}
% mdparagraph:P019:end

% mdtable:T07:begin
\par\begingroup
\fontsize{10.5}{13.2}\selectfont
\setlength{\LTleft}{0pt}\setlength{\LTright}{\fill}
\noindent
\begin{tabular}{@{}>{\raggedright\arraybackslash}p{32.31mm}>{\raggedright\arraybackslash}p{44.37mm}>{\raggedright\arraybackslash}p{41.69mm}@{}}
\toprule
\textbf{Person and number} & \textbf{Old English} & \textbf{English} \\
\midrule
First singular & Ic hæbbe ðæt hūs geworht. & I have built the house. \\
Second singular & Ðū hæfst ðæt hūs geworht. & You have built the house. \\
Third singular & Hē hæfð ðæt hūs geworht. & He has built the house. \\
First plural & Wē habbað ðæt hūs geworht. & We have built the house. \\
Second plural & Gē habbað ðæt hūs geworht. & You have built the house. \\
Third plural & Hīe habbað ðæt hūs geworht. & They have built the house. \\
\bottomrule
\end{tabular}
\par\endgroup

% mdtable:T07:end
\reserve{11\baselineskip}
% mdparagraph:P020:begin
\textbf{Preterite indicative construction.}
% mdparagraph:P020:end

% mdtable:T08:begin
\par\begingroup
\fontsize{10.5}{13.2}\selectfont
\setlength{\LTleft}{0pt}\setlength{\LTright}{\fill}
\noindent
\begin{tabular}{@{}>{\raggedright\arraybackslash}p{32.64mm}>{\raggedright\arraybackslash}p{44.83mm}>{\raggedright\arraybackslash}p{40.9mm}@{}}
\toprule
\textbf{Person and number} & \textbf{Old English} & \textbf{English} \\
\midrule
First singular & Ic hæfde ðæt hūs geworht. & I had built the house. \\
Second singular & Ðū hæfdest ðæt hūs geworht. & You had built the house. \\
Third singular & Hē hæfde ðæt hūs geworht. & He had built the house. \\
First plural & Wē hæfdon ðæt hūs geworht. & We had built the house. \\
Second plural & Gē hæfdon ðæt hūs geworht. & You had built the house. \\
Third plural & Hīe hæfdon ðæt hūs geworht. & They had built the house. \\
\bottomrule
\end{tabular}
\par\endgroup

% mdtable:T08:end
% mdparagraph:P021:begin
\textbf{Explanation:} The finite auxiliaries are \textbf{hæbbe, hæfst, hæfð, habbað} and \textbf{hæfde, hæfdest, hæfdon}, assigned to the labeled persons and numbers. \textbf{Ðæt hūs} remains the accusative object; \textbf{geworht} remains the selected past participle. The present construction presents completed building from a present viewpoint; the past construction places it before a past viewpoint. These are twelve construction positions, not twelve additional synthetic forms. Other grammatical moods or different objects do not supply the requested positions. The unchanged base does not mean that every attested participle must be uninflected: valid object agreement is distinguished from the finite auxiliary's subject agreement. [3]
% mdparagraph:P021:end

\reserve{8\baselineskip}
\par\begin{minipage}{\linewidth}\raggedright
% mdparagraph:P022:begin
\textbf{F07. Present and past participle bases.}
% mdparagraph:P022:end

% item: NEW-U12-F07
% key: KEY-NEW-U12-F07
% mdtable:T09:begin
\par\begingroup
\fontsize{10.5}{13.2}\selectfont
\setlength{\LTleft}{0pt}\setlength{\LTright}{\fill}
\noindent
\begin{tabular}{@{}>{\raggedright\arraybackslash}p{33.02mm}>{\raggedright\arraybackslash}p{44.91mm}>{\raggedright\arraybackslash}p{40.45mm}@{}}
\toprule
\textbf{Verb} & \textbf{Present participle} & \textbf{Past participle} \\
\midrule
gierwan & gierwende & gierwed \\
sęttan & sęttende & sętt \\
gewyrcean & gewyrcende & geworht \\
\bottomrule
\end{tabular}
\par\endgroup

% mdtable:T09:end
% mdparagraph:P023:begin
\textbf{Explanation:} The \textbf{past} bases \textbf{gierwed, sętt, geworht} enter the selected completed-action construction with \textbf{habban}. The present bases in \textbf{-ende} do not replace them there. Each pair is nonfinite; it supplies no person ending and does not itself assign a complete participial adjective declension. Preserve the lexical stem change in \textbf{gewyrcende/\allowbreak{}geworht} and the contracted \textbf{sętt}. The documented participle variants may be accepted with their correct analysis; a finite preterite in the participle position would change the construction. [5]
% mdparagraph:P023:end

\end{minipage}\par\addvspace{4pt}

\reserve{7\baselineskip}
% mdparagraph:P024:begin
\section*{A. Tense, sequence and voice}
% mdparagraph:P024:end

\reserve{5\baselineskip}
\par\begin{minipage}{\linewidth}\raggedright
% mdparagraph:P025:begin
\textbf{A01. Successive events.}
% mdparagraph:P025:end

% item: NEW-U12-A01
% key: KEY-NEW-U12-A01
% mdparagraph:P026:begin
\textbf{Se mann cōm, and hē geseah ðæt hūs.}
% mdparagraph:P026:end

% mdparagraph:P027:begin
The man came, and he saw the house.
% mdparagraph:P027:end

% mdparagraph:P028:begin
Arrival precedes seeing in this ordinary narrative sequence. \textbf{Cōm} and \textbf{geseah} are third-singular preterite indicatives. Equivalent English past narration is acceptable; neither verb is an infinitive, and the coordination does not state that one event caused the other.
% mdparagraph:P028:end

\end{minipage}\par\addvspace{4pt}

\par\begin{minipage}{\linewidth}\raggedright
% mdparagraph:P029:begin
\textbf{A02. An earlier event in a simple preterite.}
% mdparagraph:P029:end

% item: NEW-U12-A02
% key: KEY-NEW-U12-A02
% mdparagraph:P030:begin
\textbf{Context:} The house was finished before the visit.
% mdparagraph:P030:end

% mdparagraph:P031:begin
\textbf{Se mann geseah ðæt hūs ðe his sunu geworhte.}
% mdparagraph:P031:end

% mdparagraph:P032:begin
The man saw the house which his son had built.
% mdparagraph:P032:end

% mdparagraph:P033:begin
Both \textbf{geseah} and \textbf{geworhte} are morphologically preterite. The context makes the son's building earlier than the man's seeing; English \textbf{had built} makes that relation explicit. \textbf{His sunu} is the builder, and relative \textbf{ðe} represents the house as object of \textbf{geworhte}. \textbf{Which his son built} is also possible English if the supplied sequence remains clear. Do not call \textbf{geworhte} a morphological pluperfect or place the building after the visit. [2]
% mdparagraph:P033:end

\end{minipage}\par\addvspace{4pt}

\par\begin{minipage}{\linewidth}\raggedright
% mdparagraph:P034:begin
\textbf{A03. Present and past viewpoints.}
% mdparagraph:P034:end

% item: NEW-U12-A03
% key: KEY-NEW-U12-A03
% mdparagraph:P035:begin
\textbf{Context:} The speaker describes completed building work that they performed.
% mdparagraph:P035:end

% mdparagraph:P036:begin
\textbf{Ic hæbbe ðæt hūs geworht.} — I have built the house.
% mdparagraph:P036:end

% mdparagraph:P037:begin
\textbf{Ic hæfde ðæt hūs geworht.} — I had built the house.
% mdparagraph:P037:end

% mdparagraph:P038:begin
\textbf{Ic} is the subject and builder; \textbf{ðæt hūs} is the accusative object. \textbf{Hæbbe/\allowbreak{}hæfde} is the finite auxiliary; \textbf{geworht} is the past participle. The first clause presents completed building from the present; the second places it before a past viewpoint. The participle stays nonfinite in both. Natural contextual English alternatives are acceptable if they preserve the given speaker's agency and the different viewpoints. Neither sentence makes the house the grammatical subject or leaves the builder unnamed. [3]
% mdparagraph:P038:end

\end{minipage}\par\addvspace{4pt}

\par\begin{minipage}{\linewidth}\raggedright
% mdparagraph:P039:begin
\textbf{A04. Possession and completed action.}
% mdparagraph:P039:end

% item: NEW-U12-A04
% key: KEY-NEW-U12-A04
% mdparagraph:P040:begin
\textbf{Ic hæbbe ðæt hūs.} — I have the house.
% mdparagraph:P040:end

% mdparagraph:P041:begin
\textbf{Context for the second sentence:} The speaker has just finished building the house.
% mdparagraph:P041:end

% mdparagraph:P042:begin
\textbf{Ic hæbbe ðæt hūs geworht.} — I have built the house.
% mdparagraph:P042:end

% mdparagraph:P043:begin
In the first sentence \textbf{habban} expresses possession. In the second it combines with \textbf{geworht} to present the speaker's completed building and its present result. \textbf{Ðæt hūs} remains accusative in both. Possession alone does not assert that the possessor built the house, and \textbf{hæbbe} alone does not mean \textbf{build}. A resultative English paraphrase is defensible if it preserves the instructed builder and completion; an unsupported causative interpretation would change the supplied situation.
% mdparagraph:P043:end

\end{minipage}\par\addvspace{4pt}

\par\begin{minipage}{\linewidth}\raggedright
% mdparagraph:P044:begin
\textbf{A05. Active and passive roles.}
% mdparagraph:P044:end

% item: NEW-U12-A05
% key: KEY-NEW-U12-A05
% mdtable:T10:begin
\par\begingroup
\fontsize{10.5}{13.2}\selectfont
\setlength{\LTleft}{0pt}\setlength{\LTright}{\fill}
\noindent
\begin{tabular}{@{}>{\raggedright\arraybackslash}p{54.46mm}>{\raggedright\arraybackslash}p{66.73mm}@{}}
\toprule
\textbf{Old English} & \textbf{Translation} \\
\midrule
Se mann band ðæt hors. & The man bound the horse. \\
Ðæt hors wearð bunden. & The horse was bound. \\
Ðæt hors wæs bunden. & The horse was bound. \\
\bottomrule
\end{tabular}
\par\endgroup

% mdtable:T10:end
% mdparagraph:P045:begin
In the active clause \textbf{se mann} is the nominative subject and \textbf{ðæt hors} the accusative object. In both passives \textbf{ðæt hors} is nominative subject. The horse is the affected participant throughout; only the active sentence names the binder. Its neuter nominative and accusative spellings happen to coincide.
% mdparagraph:P045:end

% mdparagraph:P046:begin
\textbf{Wearð} commonly foregrounds the binding event; \textbf{wæs} commonly presents the resulting state. \textbf{Became bound} and \textbf{was in a bound state} can make those tendencies explicit, but context also permits event readings with \textbf{wesan}. The forms alone do not impose two different English tenses. Do not invent an agent in either passive or keep the active agent as the passive subject. [2]
% mdparagraph:P046:end

\end{minipage}\par\addvspace{4pt}

\par\begin{minipage}{\linewidth}\raggedright
% mdparagraph:P047:begin
\textbf{A06. An intransitive past participle.}
% mdparagraph:P047:end

% item: NEW-U12-A06
% key: KEY-NEW-U12-A06
% mdparagraph:P048:begin
\textbf{Se mann wæs gewiten.}
% mdparagraph:P048:end

% mdparagraph:P049:begin
The man had departed; or, the man was gone.
% mdparagraph:P049:end

% mdparagraph:P050:begin
\textbf{Wæs} is finite third-singular preterite \textbf{wesan}; \textbf{gewiten} is the past participle of intransitive \textbf{gewītan}. The man himself has departed. He is not the object of an unnamed agent's transitive action. \textbf{Had departed} makes anteriority explicit; \textbf{was gone} emphasizes the resulting absence. This recognition example does not require an invented passive agent or establish a general rule for all English \textbf{be + participle} expressions. [2]
% mdparagraph:P050:end

\end{minipage}\par\addvspace{4pt}

\par\begin{minipage}{\linewidth}\raggedright
% mdparagraph:P051:begin
\textbf{A07. A present participle in a past setting.}
% mdparagraph:P051:end

% item: NEW-U12-A07
% key: KEY-NEW-U12-A07
% mdparagraph:P052:begin
\textbf{Se mann wæs sittende.}
% mdparagraph:P052:end

% mdparagraph:P053:begin
The man sat; or, the man was sitting.
% mdparagraph:P053:end

% mdparagraph:P054:begin
\textbf{Wæs} is the finite third-singular preterite; \textbf{sittende} is a present participle. An ongoing situation can make \textbf{was sitting} idiomatic English, but the Old English expression does not mechanically require the modern English progressive. \textbf{Sat} is a valid simple-past rendering. The name \textbf{present participle} identifies the nonfinite form, not the time of the whole sentence; \textbf{sittende} is neither a finite present nor a past participle. [2]
% mdparagraph:P054:end

\end{minipage}\par\addvspace{4pt}

\par\begin{minipage}{\linewidth}\raggedright
% mdparagraph:P055:begin
\textbf{A08. Presence simultaneous with a report.}
% mdparagraph:P055:end

% item: NEW-U12-A08
% key: KEY-NEW-U12-A08
% mdparagraph:P056:begin
\textbf{Context:} In each sentence the presence is simultaneous with the saying.
% mdparagraph:P056:end

% mdparagraph:P057:begin
\textbf{Hē sægð ðæt hē sīe hēr.} — He says that he is here.
% mdparagraph:P057:end

% mdparagraph:P058:begin
\textbf{Hē sægde ðæt hē wǣre hēr.} — He said that he was here.
% mdparagraph:P058:end

% mdparagraph:P059:begin
\textbf{Sægð} is present indicative; \textbf{sægde} is preterite indicative in the main narrative clause. Dependent \textbf{sīe} is present subjunctive and \textbf{wǣre} preterite subjunctive. The first presence is related to the present saying; the second to the past saying. The second sentence does not, in this context, mean that the presence ended before he spoke. The report's subjunctive does not assert absence or falsity. Equivalent indirect English reports are valid; do not infer \textbf{had been} merely from preterite \textbf{wǣre}. [2]
% mdparagraph:P059:end

\end{minipage}\par\addvspace{4pt}

\par\begin{minipage}{\linewidth}\raggedright
% mdparagraph:P060:begin
\textbf{A09. Building anterior to a report.}
% mdparagraph:P060:end

% item: NEW-U12-A09
% key: KEY-NEW-U12-A09
% mdparagraph:P061:begin
\textbf{Context:} The building precedes the saying.
% mdparagraph:P061:end

% mdparagraph:P062:begin
\textbf{Hē sægde ðæt hē ðæt hūs geworht hæfde.}
% mdparagraph:P062:end

% mdparagraph:P063:begin
He said that he had built the house.
% mdparagraph:P063:end

% mdparagraph:P064:begin
Main-clause \textbf{sægde} is the finite preterite indicative. Dependent \textbf{hæfde} is the finite preterite subjunctive in the assigned reported-statement analysis; \textbf{geworht} is the past participle. \textbf{Hē} is the reported builder and \textbf{ðæt hūs} the accusative object. The auxiliary–participle construction presents completed building before the past report. \textbf{Hæfde} has the same spelling as indicative forms, so its spelling alone does not identify mood. A grammatical alternative order retains the same roles and relation; \textbf{geworht} must not be analyzed as another finite verb. [2–3]
% mdparagraph:P064:end

\end{minipage}\par\addvspace{4pt}

\par\begin{minipage}{\linewidth}\raggedright
% mdparagraph:P065:begin
\textbf{A10. Reported intention.}
% mdparagraph:P065:end

% item: NEW-U12-A10
% key: KEY-NEW-U12-A10
% mdparagraph:P066:begin
\textbf{Hē sægde ðæt hē wolde cuman.}
% mdparagraph:P066:end

% mdparagraph:P067:begin
He said that he intended to come; or, that he wished/\allowbreak{}wanted to come.
% mdparagraph:P067:end

% mdparagraph:P068:begin
The intended action is the coming, expressed by infinitive \textbf{cuman} with \textbf{wolde}. It is prospective from the past saying. The sentence reports the intention without asserting an actual arrival. \textbf{Would come} can convey that intention in suitable English context, but does not justify an automatic future-tense analysis. \textbf{He came} would assert an event which this sentence leaves open. [2]
% mdparagraph:P068:end

\end{minipage}\par\addvspace{4pt}

\par\begin{minipage}{\linewidth}\raggedright
% mdparagraph:P069:begin
\textbf{A11. Expected arrival and desire.}
% mdparagraph:P069:end

% item: NEW-U12-A11
% key: KEY-NEW-U12-A11
% mdparagraph:P070:begin
\textbf{Context for the first sentence:} The visit is still expected.
% mdparagraph:P070:end

% mdparagraph:P071:begin
\textbf{Hē cymð.} — He will come; or, he is coming.
% mdparagraph:P071:end

% mdparagraph:P072:begin
\textbf{Hē wille cuman.} — He wishes/\allowbreak{}intends to come.
% mdparagraph:P072:end

% mdparagraph:P073:begin
In the first, \textbf{cymð} is morphologically present indicative, with future reference supplied by the situation. In the second, \textbf{wille} contributes desire or intention and \textbf{cuman} names the intended action. \textbf{Wants to come} is also apt. English \textbf{will} therefore need not be translated with \textbf{willan}; conversely, the stated volition in the second sentence should not disappear into a purely mechanical tense label.
% mdparagraph:P073:end

\end{minipage}\par\addvspace{4pt}

\reserve{7\baselineskip}
% mdparagraph:P074:begin
\section*{C. Controlled composition}
% mdparagraph:P074:end

\reserve{5\baselineskip}
\par\begin{minipage}{\linewidth}\raggedright
% mdparagraph:P075:begin
\textbf{C01. We have built the house.}
% mdparagraph:P075:end

% item: NEW-U12-C01
% key: KEY-NEW-U12-C01
% mdparagraph:P076:begin
\textbf{Wē habbað ðæt hūs geworht.}
% mdparagraph:P076:end

% mdparagraph:P077:begin
\textbf{Habbað} agrees with first-plural \textbf{wē}. The house remains the accusative object and \textbf{geworht} the selected past participle. Other grammatical positions of the object and participle preserve the target. Singular \textbf{hæbbe} would not agree with \textbf{wē}; a simple preterite might be appropriate in another context but would omit the requested present \textbf{habban} construction.
% mdparagraph:P077:end

\end{minipage}\par\addvspace{4pt}

\par\begin{minipage}{\linewidth}\raggedright
% mdparagraph:P078:begin
\textbf{C02. They had built the house.}
% mdparagraph:P078:end

% item: NEW-U12-C02
% key: KEY-NEW-U12-C02
% mdparagraph:P079:begin
\textbf{Hīe hæfdon ðæt hūs geworht.}
% mdparagraph:P079:end

% mdparagraph:P080:begin
The preterite plural auxiliary \textbf{hæfdon} puts the completed building before the supplied past viewpoint. Other grammatical word orders are acceptable. A simple preterite can express an earlier event in context, as A02 shows, but does not meet this task's explicit compound-construction target. Present \textbf{habbað} would change the requested viewpoint.
% mdparagraph:P080:end

\end{minipage}\par\addvspace{4pt}

\par\begin{minipage}{\linewidth}\raggedright
% mdparagraph:P081:begin
\textbf{C03. The men and the horse.}
% mdparagraph:P081:end

% item: NEW-U12-C03
% key: KEY-NEW-U12-C03
% mdparagraph:P082:begin
\textbf{Ðā męnn bundon ðæt hors.} — The men bound the horse.
% mdparagraph:P082:end

% mdparagraph:P083:begin
\textbf{Ðæt hors wearð bunden.} — The horse was bound.
% mdparagraph:P083:end

% mdparagraph:P084:begin
The plural active subject requires \textbf{bundon}. The passive makes singular \textbf{hors} the subject and therefore has singular \textbf{wearð}. It leaves the binder unexpressed, so the information that men performed the action is no longer stated. A passive with \textbf{wæs} can be grammatical but misses the requested \textbf{weorðan} target. Plural \textbf{wurdon} would wrongly retain agreement with the former active subject.
% mdparagraph:P084:end

\end{minipage}\par\addvspace{4pt}

\par\begin{minipage}{\linewidth}\raggedright
% mdparagraph:P085:begin
\textbf{C04. The monk's reported presence.}
% mdparagraph:P085:end

% item: NEW-U12-C04
% key: KEY-NEW-U12-C04
% mdparagraph:P086:begin
\textbf{Se munuc sægð ðæt hē sīe hēr.}
% mdparagraph:P086:end

% mdparagraph:P087:begin
\textbf{Se munuc sægde ðæt hē wǣre hēr.}
% mdparagraph:P087:end

% mdparagraph:P088:begin
The monk's presence is simultaneous with the saying in each version. \textbf{Sīe/\allowbreak{}wǣre} supplies the requested report subjunctive. An indicative report may be valid Old English in an appropriate presentation, but it does not satisfy that stated choice. Other documented copular subjunctive forms or grammatical word orders are acceptable with the same temporal relation. Do not change simultaneous presence to an earlier \textbf{had been} or make \textbf{hē} refer to a different person without context.
% mdparagraph:P088:end

\end{minipage}\par\addvspace{4pt}

\par\begin{minipage}{\linewidth}\raggedright
% mdparagraph:P089:begin
\textbf{C05. Building before the report.}
% mdparagraph:P089:end

% item: NEW-U12-C05
% key: KEY-NEW-U12-C05
% mdparagraph:P090:begin
\textbf{Hē sægde ðæt hē ðæt hūs geworht hæfde.}
% mdparagraph:P090:end

% mdparagraph:P091:begin
The dependent finite form is preterite subjunctive \textbf{hæfde}; the past participle is \textbf{geworht}. The completed action precedes the saying. \textbf{Hē sægde ðæt hē hæfde ðæt hūs geworht} is an acceptable alternative order. The identical indicative spelling of \textbf{hæfde} does not remove the task's specified mood analysis. A passive with the house as subject would change the participant structure. [2–3]
% mdparagraph:P091:end

\end{minipage}\par\addvspace{4pt}

\par\begin{minipage}{\linewidth}\raggedright
% mdparagraph:P092:begin
\textbf{C06. An intention to come.}
% mdparagraph:P092:end

% item: NEW-U12-C06
% key: KEY-NEW-U12-C06
% mdparagraph:P093:begin
\textbf{Hē sægde ðæt hē wolde cuman.}
% mdparagraph:P093:end

% mdparagraph:P094:begin
\textbf{Wolde cuman} expresses the intention from the past reporting situation; the infinitive is not an assertion that the coming occurred. Defensible alternative word order is acceptable. Replacing the dependent expression with an asserted past arrival would lose the intended meaning, while translating \textbf{wolde} as \textbf{intended}, \textbf{wished} or \textbf{wanted} preserves it.
% mdparagraph:P094:end

\end{minipage}\par\addvspace{4pt}

\par\begin{minipage}{\linewidth}\raggedright
% mdparagraph:P095:begin
\textbf{C07. Preparing food.}
% mdparagraph:P095:end

% item: NEW-U12-C07
% key: KEY-NEW-U12-C07
% mdparagraph:P096:begin
\textbf{Giere ðone męte.} — Prepare the food.
% mdparagraph:P096:end

% mdparagraph:P097:begin
\textbf{Ic hæbbe ðone męte gierwed.} — I have prepared the food.
% mdparagraph:P097:end

% mdparagraph:P098:begin
\textbf{Giere} addresses one person. The food is the accusative object in both clauses. In the reply \textbf{hæbbe} is finite and \textbf{gierwed} the past participle. Documented forms of \textbf{gierwan}, correctly analyzed object-agreeing participles and grammatical word-order alternatives can be accepted. Plural \textbf{gierwað} does not meet the singular address; a finite preterite substituted for \textbf{gierwed} would not complete the compound construction.
% mdparagraph:P098:end

\end{minipage}\par\addvspace{4pt}

\par\begin{minipage}{\linewidth}\raggedright
% mdparagraph:P099:begin
\textbf{C08. The completed bridge and a new arrival.}
% mdparagraph:P099:end

% item: NEW-U12-C08
% key: KEY-NEW-U12-C08
% mdparagraph:P100:begin
\textbf{Sēo brycg wæs geworht. Ðā cōm mann.}
% mdparagraph:P100:end

% mdparagraph:P101:begin
The bridge is the feminine subject of the first clause; \textbf{wæs geworht} presents its completed state without naming a builder. The second clause introduces a man with finite preterite \textbf{cōm}. Indefinite English \textbf{a} does not require cardinal \textbf{ān}. Other grammatical word orders are possible. A suitable passive with \textbf{weorðan} can be valid language but misses the specified \textbf{wesan} here; a masculine or neuter article with \textbf{brycg} would violate agreement.
% mdparagraph:P101:end

\end{minipage}\par\addvspace{4pt}

\reserve{7\baselineskip}
% mdparagraph:P102:begin
\section*{Cook, Group VI, Exercise 19}
% mdparagraph:P102:end

\par\begin{minipage}{\linewidth}\raggedright
% mdparagraph:P103:begin
\textbf{19.1.} At length the king ceased to bear rule.
% mdparagraph:P103:end

% item: EX-VI-19-01
% key: KEY-EX-VI-19-01
% mdparagraph:P104:begin
\textbf{Æt nīehstan se cyning geswāc rīces.}
% mdparagraph:P104:end

% mdparagraph:P105:begin
\textbf{Geswāc} is the simple preterite of \textbf{geswīcan}; \textbf{rīces} is the genitive complement, rule or dominion. The supplied \textbf{æt nīehstan} gives \textbf{at length}. The construction conveys the king's cessation of rule without a separate infinitive meaning \textbf{bear}. Equivalent supported wording is possible, but an invented infinitival frame would not follow the taught government.
% mdparagraph:P105:end

\end{minipage}\par\addvspace{4pt}

\par\begin{minipage}{\linewidth}\raggedright
% mdparagraph:P106:begin
\textbf{19.2.} He had the kingdom forty years.
% mdparagraph:P106:end

% item: EX-VI-19-02
% key: KEY-EX-VI-19-02
% mdparagraph:P107:begin
\textbf{Hē hæfde ðæt rīce fēowertig gēara.}
% mdparagraph:P107:end

% mdparagraph:P108:begin
Here \textbf{hæfde} means possessed or held: it is lexical \textbf{habban}, not a perfect auxiliary. \textbf{Ðæt rīce} is its accusative object. The supplied \textbf{fēowertig gēara} expresses the forty-year duration; \textbf{gēara} is the genitive plural used with the numeral. Do not add a participle merely because English has \textbf{had}, or turn the kingdom into the subject. \textbf{He held the kingdom for forty years} gives the same meaning.
% mdparagraph:P108:end

\end{minipage}\par\addvspace{4pt}

\par\begin{minipage}{\linewidth}\raggedright
% mdparagraph:P109:begin
\textbf{19.3.} When were the bridges built?
% mdparagraph:P109:end

% item: EX-VI-19-03
% key: KEY-EX-VI-19-03
% mdparagraph:P110:begin
\textbf{Hwonne wurdon ðā brycga geworht?}
% mdparagraph:P110:end

% mdparagraph:P111:begin
\textbf{Hwonne} asks when; \textbf{ðā brycga} is the plural subject and \textbf{wurdon} the plural preterite auxiliary. The model foregrounds the building event. \textbf{Wǣron … geworht} can also be defensible when context supports that passive presentation; \textbf{wesan} is not restricted to states in every occurrence. The selected bare participle is legitimate with a plural subject. Singular \textbf{wearð} would fail to agree, and an active answer with a newly invented builder would add information.
% mdparagraph:P111:end

\end{minipage}\par\addvspace{4pt}

\par\begin{minipage}{\linewidth}\raggedright
% mdparagraph:P112:begin
\textbf{19.4.} To-day thou shalt be with me.
% mdparagraph:P112:end

% item: EX-VI-19-04
% key: KEY-EX-VI-19-04
% mdparagraph:P113:begin
\textbf{Tō-dæg ðū bist mid mē.}
% mdparagraph:P113:end

% mdparagraph:P114:begin
The present form \textbf{bist} has future reference in this promise; \textbf{mid mē} names the accompanying person in the dative. The sentence does not need \textbf{willan} simply because the English has \textbf{shalt}. \textbf{Today you will be with me} is the intended meaning. A translation introducing personal desire or an obligation to be present would add a different modal meaning.
% mdparagraph:P114:end

\end{minipage}\par\addvspace{4pt}

\par\begin{minipage}{\linewidth}\raggedright
% mdparagraph:P115:begin
\textbf{19.5.} What (om. did) the Romans build (pret.) in Britain?
% mdparagraph:P115:end

% item: EX-VI-19-05
% key: KEY-EX-VI-19-05
% mdparagraph:P116:begin
\textbf{Hwæt geworhton Rōmane on Breotone?}
% mdparagraph:P116:end

% mdparagraph:P117:begin
\textbf{Hwæt} is the thing built, and \textbf{Rōmane} the supplied plural subject. \textbf{Geworhton} is the explicitly requested preterite; no separate word translates auxiliary English \textbf{did}. The local \textbf{on Breotone} supplies the location. A grammatical order with the subject earlier is acceptable. Do not put the question word into the subject role or use a present verb against Cook's \textbf{(pret.)} control.
% mdparagraph:P117:end

\end{minipage}\par\addvspace{4pt}

\par\begin{minipage}{\linewidth}\raggedright
% mdparagraph:P118:begin
\textbf{19.6.} I send my angel before (beforan, 166) thy face.
% mdparagraph:P118:end

% item: EX-VI-19-06
% key: KEY-EX-VI-19-06
% mdparagraph:P119:begin
\textbf{Ic sęnde mīnne ęngel beforan ðīnum andwlitan.}
% mdparagraph:P119:end

% mdparagraph:P120:begin
\textbf{Mīnne ęngel} is the masculine accusative person sent. \textbf{Beforan} takes the taught dative phrase \textbf{ðīnum andwlitan} here, with the possessive agreeing with \textbf{andwlitan}. The literal reference to the addressee's face is retained. Grammatical word-order alternatives are acceptable; using an accusative after this selected \textbf{beforan} frame, or making the angel dative without a new construction, would change the required relations.
% mdparagraph:P120:end

\end{minipage}\par\addvspace{4pt}

\par\begin{minipage}{\linewidth}\raggedright
% mdparagraph:P121:begin
\textbf{19.7.} Come (sing.), and bring those letters that you wrote.
% mdparagraph:P121:end

% item: EX-VI-19-07
% key: KEY-EX-VI-19-07
% mdparagraph:P122:begin
\textbf{Cum, and bring ðā gewritu ðe ðū write.}
% mdparagraph:P122:end

% mdparagraph:P123:begin
\textbf{Cum} and \textbf{bring} address one person, as Cook expressly requires for \textbf{come}. \textbf{Ðā gewritu} is the plural object, and \textbf{ðe} represents it as object of \textbf{write}. The unit's directions also select singular \textbf{ðū} in the relative clause: preterite second-singular \textbf{write} has short \textbf{i}, unlike present first-singular \textbf{wrīte}. Cook's unnumbered relative \textbf{you} could instead yield grammatical \textbf{ðe gē writon} under a plural interpretation, but that does not meet this unit's selected singular target. Do not change the first command to plural or silently read \textbf{write} as a present.
% mdparagraph:P123:end

\end{minipage}\par\addvspace{4pt}

\par\begin{minipage}{\linewidth}\raggedright
% mdparagraph:P124:begin
\textbf{19.8.} I promise you obedience, if you help (196. d) me.
% mdparagraph:P124:end

% item: EX-VI-19-08
% key: KEY-EX-VI-19-08
% mdparagraph:P125:begin
\textbf{Ic gehāte ēow gehīersumnesse, gif gē mē helpen.}
% mdparagraph:P125:end

% mdparagraph:P126:begin
The model takes both occurrences of \textbf{you} as plural: \textbf{ēow} is the dative promisee, \textbf{gehīersumnesse} the promised thing, and \textbf{gē … helpen} the conditional clause. The reference \textbf{196. d} selects the subjunctive \textbf{helpen}; \textbf{mē} is the dative person helped. Singular \textbf{ðē/\allowbreak{}ðū … helpe} is also possible if the addressee is consistently singular. An indicative can occur in other conditional presentations, but does not satisfy this source control. Do not give the recipient the case of the promised thing.
% mdparagraph:P126:end

\end{minipage}\par\addvspace{4pt}

\par\begin{minipage}{\linewidth}\raggedright
% mdparagraph:P127:begin
\textbf{19.9.} Send me (om. a) large army.
% mdparagraph:P127:end

% item: EX-VI-19-09
% key: KEY-EX-VI-19-09
% mdparagraph:P128:begin
\textbf{Sęnd mē micelne hęre.}
% mdparagraph:P128:end

% mdparagraph:P129:begin
\textbf{Sęnd} is the singular command selected by the unit's directions; \textbf{mē} is the dative recipient and \textbf{micelne hęre} the masculine accusative object. Cook's English does not specify the addressee's number, so plural \textbf{sęndað} is grammatical under another interpretation, but it misses the unit's singular target. The adjective's \textbf{-ne} marks the object's agreement even though \textbf{hęre} itself has the same nominative and accusative spelling. Omit an indefinite article as directed; cardinal \textbf{ān} is not required.
% mdparagraph:P129:end

\end{minipage}\par\addvspace{4pt}

\par\begin{minipage}{\linewidth}\raggedright
% mdparagraph:P130:begin
\textbf{19.10.} I fight with my enemies.
% mdparagraph:P130:end

% item: EX-VI-19-10
% key: KEY-EX-VI-19-10
% mdparagraph:P131:begin
\textbf{Ic feohte wið mīne fīend.}
% mdparagraph:P131:end

% mdparagraph:P132:begin
\textbf{Feohte} is first-singular present indicative. The selected opponent frame uses \textbf{wið + accusative}: \textbf{mīne fīend}, my enemies. The possessive agrees with the plural noun. This \textbf{with} identifies opponents, not companions. Other evidenced government may be valid in its own context but would miss this selected accusative target. The source reading's separately supplied \textbf{wið + dative} must not be treated as the only possible construction here.
% mdparagraph:P132:end

\end{minipage}\par\addvspace{4pt}

\par\begin{minipage}{\linewidth}\raggedright
% mdparagraph:P133:begin
\textbf{19.11.} Expel your foes.
% mdparagraph:P133:end

% item: EX-VI-19-11
% key: KEY-EX-VI-19-11
% mdparagraph:P134:begin
\textbf{Ādrīfað ēowre fīend.}
% mdparagraph:P134:end

% mdparagraph:P135:begin
The unit's directions select several addressees: plural \textbf{ādrīfað} matches possessive \textbf{ēowre}. \textbf{Fīend} is the accusative plural object. Singular \textbf{Ādrīf ðīne fīend} is grammatical for another interpretation of Cook's unnumbered English, but does not meet the unit's plural target. Keep command number and the possessor consistent. A finite indicative substituted for the command, or a possessive that changes the selected addressee, would miss the target.
% mdparagraph:P135:end

\end{minipage}\par\addvspace{4pt}

\par\begin{minipage}{\linewidth}\raggedright
% mdparagraph:P136:begin
\textbf{19.12.} The soldiers went home.
% mdparagraph:P136:end

% item: EX-VI-19-12
% key: KEY-EX-VI-19-12
% mdparagraph:P137:begin
\textbf{Ðā cęmpan fōron hām.}
% mdparagraph:P137:end

% mdparagraph:P138:begin
Plural \textbf{ðā cęmpan} requires preterite plural \textbf{fōron}. Directional \textbf{hām} means home and needs no preposition. Another taught verb of going can be defensible if it preserves the meaning and reference. Do not use singular \textbf{fōr} with the soldiers or turn \textbf{hām} into a static location by adding a different construction.
% mdparagraph:P138:end

\end{minipage}\par\addvspace{4pt}

\par\begin{minipage}{\linewidth}\raggedright
% mdparagraph:P139:begin
\textbf{19.13.} He will come with (om. a) fleet.
% mdparagraph:P139:end

% item: EX-VI-19-13
% key: KEY-EX-VI-19-13
% mdparagraph:P140:begin
\textbf{Hē cymð mid sciphęre.}
% mdparagraph:P140:end

% mdparagraph:P141:begin
Present \textbf{cymð} has future reference in the supplied English context. \textbf{Mid sciphęre} expresses accompaniment in the dative; the collective fleet noun does not require a plural simply because it includes several ships. Cook expressly omits \textbf{a}. Adding \textbf{willan} would introduce volition unless a suitable context justified it, rather than merely supplying a missing future-tense marker.
% mdparagraph:P141:end

\end{minipage}\par\addvspace{4pt}

\par\begin{minipage}{\linewidth}\raggedright
% mdparagraph:P142:begin
\textbf{19.14.} Slay (sing.) all that you find (gemētan).
% mdparagraph:P142:end

% item: EX-VI-19-14
% key: KEY-EX-VI-19-14
% mdparagraph:P143:begin
\textbf{Sleah ealle ðā ðe ðū gemētst.}
% mdparagraph:P143:end

% mdparagraph:P144:begin
\textbf{Sleah} is the singular command, with the newly assigned lethal sense. \textbf{Ealle ðā ðe} gives all those whom, with the people as object of \textbf{slēan} and as the understood relative object of \textbf{gemētst}. \textbf{Ðū gemētst} retains the singular addressee. A different grammatical word order can preserve these relations. Do not use plural \textbf{slēað}, change the victims into finders, or translate the assigned sense merely as a nonlethal strike.
% mdparagraph:P144:end

\end{minipage}\par\addvspace{4pt}

\par\begin{minipage}{\linewidth}\raggedright
% mdparagraph:P145:begin
\textbf{19.15.} The harvest is rich in fruits.
% mdparagraph:P145:end

% item: EX-VI-19-15
% key: KEY-EX-VI-19-15
% mdparagraph:P146:begin
\textbf{Se hærfest is welig on wæstmum.}
% mdparagraph:P146:end

% mdparagraph:P147:begin
\textbf{Welig} is the masculine nominative predicate agreeing with \textbf{hærfest}. The rich-in frame uses \textbf{on + dative}, here plural \textbf{wæstmum}. Preserve short \textbf{æ} in the checked noun \textbf{wæstm}. Equivalent word order is acceptable; changing the phrase to an accusative destination would lose the respect in which the harvest is rich.
% mdparagraph:P147:end

\end{minipage}\par\addvspace{4pt}

\par\begin{minipage}{\linewidth}\raggedright
% mdparagraph:P148:begin
\textbf{19.16.} Seedtime and harvest never shall fail (geswīcan).
% mdparagraph:P148:end

% item: EX-VI-19-16
% key: KEY-EX-VI-19-16
% mdparagraph:P149:begin
\textbf{Sǣdtīma and hærfest ne geswīcað nǣfre.}
% mdparagraph:P149:end

% mdparagraph:P150:begin
The two coordinated subjects take plural \textbf{geswīcað}. This assigned intransitive use means cease or fail without an expressed genitive complement. Present tense carries the future reference. \textbf{Ne … nǣfre} is one negative-concord expression: \textbf{shall never cease}. It is not a pair of negatives canceling each other. Do not add an object simply because \textbf{geswīcan} had a genitive complement in 19.1, or treat the promise as an assertion of the seasons' desire.
% mdparagraph:P150:end

\end{minipage}\par\addvspace{4pt}

\par\begin{minipage}{\linewidth}\raggedright
% mdparagraph:P151:begin
\textbf{19.17.} The year has four seasons.
% mdparagraph:P151:end

% item: EX-VI-19-17
% key: KEY-EX-VI-19-17
% mdparagraph:P152:begin
\textbf{Ðæt gēar hæfð fēower tīda.}
% mdparagraph:P152:end

% mdparagraph:P153:begin
\textbf{Ðæt gēar} is the singular subject, with lexical \textbf{hæfð} and the accusative quantity \textbf{fēower tīda}. \textbf{Tīd} has the assigned sense season here. This is possession or inclusion, not a completed-action construction: no participle is needed. \textbf{The year contains four seasons} is equivalent English. Do not make the auxiliary plural because its object is plural.
% mdparagraph:P153:end

\end{minipage}\par\addvspace{4pt}

\reserve{7\baselineskip}
% mdparagraph:P154:begin
\section*{Cook, Group VI, Exercise 20}
% mdparagraph:P154:end

\par\begin{minipage}{\linewidth}\raggedright
% mdparagraph:P155:begin
\textbf{20.1.} The earth is (om. the) Lord’s.
% mdparagraph:P155:end

% item: EX-VI-20-01
% key: KEY-EX-VI-20-01
% mdparagraph:P156:begin
\textbf{Sēo eorðe is dryhtnes.}
% mdparagraph:P156:end

% mdparagraph:P157:begin
\textbf{Dryhtnes} is a predicative genitive: the earth belongs to the Lord. It is not an adjective agreeing with feminine \textbf{eorðe}. Cook directs omission of the article before Lord, so no separate \textbf{the} is inserted there. Another grammatical expression of possession may convey the same idea, but this task practices the newly taught genitive predicate. A dative or nominative replacement would not express that relation.
% mdparagraph:P157:end

\end{minipage}\par\addvspace{4pt}

\par\begin{minipage}{\linewidth}\raggedright
% mdparagraph:P158:begin
\textbf{20.2.} The sea saw, and it fled.
% mdparagraph:P158:end

% item: EX-VI-20-02
% key: KEY-EX-VI-20-02
% mdparagraph:P159:begin
\textbf{Se sǣ geseah, and hē flēah.}
% mdparagraph:P159:end

% mdparagraph:P160:begin
The selected \textbf{sǣ} is masculine, so \textbf{se} and referring \textbf{hē} agree with its grammatical gender. English idiom still uses \textbf{it}. \textbf{Geseah} and \textbf{flēah} are third-singular preterites, and this assigned use of seeing has no expressed object. Do not invent something seen or choose a feminine article merely because the sea is not a male person. The personification does not alter the selected declension.
% mdparagraph:P160:end

\end{minipage}\par\addvspace{4pt}

\par\begin{minipage}{\linewidth}\raggedright
% mdparagraph:P161:begin
\textbf{20.3.} He shall flee the iron weapon.
% mdparagraph:P161:end

% item: EX-VI-20-03
% key: KEY-EX-VI-20-03
% mdparagraph:P162:begin
\textbf{Hē flīehð ðæt īserne wǣpen.}
% mdparagraph:P162:end

% mdparagraph:P163:begin
The present \textbf{flīehð} has future reference. \textbf{Ðæt īserne wǣpen} is the accusative thing avoided; \textbf{īserne} is the weak neuter singular adjective after the article. This use of \textbf{flēon} does not require a phrase meaning away from. Keep the adjective, noun and article in agreement. A purely volitional paraphrase would add something to the promised or predicted flight.
% mdparagraph:P163:end

\end{minipage}\par\addvspace{4pt}

\par\begin{minipage}{\linewidth}\raggedright
% mdparagraph:P164:begin
\textbf{20.4.} He that (sē ðe) is not (om. the) shepherd, and owns not (127) the sheep, when he sees the wolf, then flees he, and abandons the sheep, and the wolf seizes the sheep.
% mdparagraph:P164:end

% item: EX-VI-20-04
% key: KEY-EX-VI-20-04
% mdparagraph:P165:begin
\textbf{Sē ðe nis scēaphierde, and nāh ðā scēap, ðonne hē gesiehð ðone wulf, ðonne flīehð hē, and forlǣtt ðā scēap, and se wulf nimð ðā scēap.}
% mdparagraph:P165:end

% mdparagraph:P166:begin
\textbf{Sē ðe} supplies the generic \textbf{he who}; \textbf{scēaphierde} has no article, as Cook directs. The local \textbf{nāh} means owns not; \textbf{nis} means is not. The first \textbf{ðonne} introduces the circumstance of seeing the wolf, and the second resumes the main action, then he flees. The nonshepherd remains the subject of \textbf{forlǣtt}; \textbf{se wulf} becomes the new subject of \textbf{nimð}. All four actions—seeing, fleeing, abandoning and seizing—must remain. Do not make the wolf abandon the sheep, omit one of the coordinated clauses, or demand an independently recalled negative paradigm from the supplied \textbf{nāh}.
% mdparagraph:P166:end

\end{minipage}\par\addvspace{4pt}

\par\begin{minipage}{\linewidth}\raggedright
% mdparagraph:P167:begin
\textbf{20.5.} The day is thine, the night also is thine; thou hast prepared the sun and the moon; thou hast set the borders of the earth; thou hast made summer and winter.
% mdparagraph:P167:end

% item: EX-VI-20-05
% key: KEY-EX-VI-20-05
% mdparagraph:P168:begin
\textbf{Se dæg is ðīn, sēo niht is ēac ðīn; ðū hæfst ðā sunnan and ðone mōnan gierwed; ðū hæfst ðā gemǣru ðǣre eorðan sętt; ðū hæfst sumer and winter geworht.}
% mdparagraph:P168:end

% mdparagraph:P169:begin
The first two clauses use predicative possessive \textbf{ðīn}, agreeing respectively with masculine \textbf{dæg} and feminine \textbf{niht}; their selected surface happens to coincide. The following three clauses use \textbf{hæfst} with accusative objects and past participles \textbf{gierwed, sętt, geworht}. \textbf{Ðǣre eorðan} depends on \textbf{gemǣru}, the borders of the earth. Preserve both heavenly bodies, the borders, and both seasons. Correctly evidenced object-agreeing participles can be valid alternatives to these selected bare bases. A simple possessive translation of all five clauses would lose the three completed actions; a participle made to agree in person with \textbf{ðū} would confuse two different kinds of agreement. [3]
% mdparagraph:P169:end

\end{minipage}\par\addvspace{4pt}

\par\begin{minipage}{\linewidth}\raggedright
% mdparagraph:P170:begin
\textbf{20.6.} Let us build these cities, and make about (ābūtan, 172) them walls, and towers, and gates, while the land is yet (gīet) before us.
% mdparagraph:P170:end

% item: EX-VI-20-06
% key: KEY-EX-VI-20-06
% mdparagraph:P171:begin
\textbf{Gewyrcen wē ðās ceastra, and gewyrcen ābūtan hīe weallas, and torras, and gatu, ðā-hwīle-ðe ðæt land is gīet beforan ūs.}
% mdparagraph:P171:end

% mdparagraph:P172:begin
Both \textbf{gewyrcen} forms are first-plural present subjunctives used as exhortations; the second continues the expressed \textbf{wē}. \textbf{Ðās ceastra} and the walls, towers and gates are accusative objects. \textbf{Ābūtan hīe} has the newly taught accusative frame; \textbf{beforan ūs} has the taught dative. \textbf{Ðā-hwīle-ðe} introduces the simultaneous condition while the land is still before them. Preserve Cook's \textbf{gīet}, with length over \textbf{i}. Repeating \textbf{gewyrcen} keeps both actions within its taught make/\allowbreak{}build sense. Do not replace it with a familiar verb whose supplied placement frame has not established this around-construction, or use indicative \textbf{gewyrcað} for the selected exhortation.
% mdparagraph:P172:end

\end{minipage}\par\addvspace{4pt}

\par\begin{minipage}{\linewidth}\raggedright
% mdparagraph:P173:begin
\textbf{20.7.} These nations are more than I; how can I expel them?
% mdparagraph:P173:end

% item: EX-VI-20-07
% key: KEY-EX-VI-20-07
% mdparagraph:P174:begin
\textbf{Ðās ðēoda sindon mā ðonne ic; hū mæg ic hīe ādrīfan?}
% mdparagraph:P174:end

% mdparagraph:P175:begin
\textbf{Mā} is the assigned indeclinable quantity predicate, more in number; it is not an inflected comparative adjective. \textbf{Ðonne ic} supplies the comparison with I, and \textbf{mæg} takes plain infinitive \textbf{ādrīfan}. \textbf{Hīe} remains the accusative nations to be expelled. An English paraphrase such as \textbf{These nations outnumber me} preserves the numerical point. Inflected \textbf{māran} would substitute a different comparative structure, and a finite expulsion verb after \textbf{mæg} would fail the modal construction.
% mdparagraph:P175:end

\end{minipage}\par\addvspace{4pt}

\par\begin{minipage}{\linewidth}\raggedright
% mdparagraph:P176:begin
\textbf{20.8.} Go to the sea, and cast a fishhook, and take the first (ǣrest) fish, and open (geopenian) his mouth (mūð, 43).
% mdparagraph:P176:end

% item: EX-VI-20-08
% key: KEY-EX-VI-20-08
% mdparagraph:P177:begin
\textbf{Gā tō ðǣm sǣ, and weorp angel, and nim ðone ǣrestan fisc, and geopena his mūð.}
% mdparagraph:P177:end

% mdparagraph:P178:begin
The four singular commands are \textbf{gā, weorp, nim, geopena}. \textbf{Tō} takes the dative; the course's selected masculine \textbf{sǣ} therefore has \textbf{ðǣm}, not feminine \textbf{ðǣre}. The fishhook is \textbf{angel}, distinct from \textbf{ęngel}, angel. Cook's bare hint \textbf{ǣrest} is realized in the taught weak ordinal \textbf{ðone ǣrestan fisc}, the first fish. \textbf{His mūð} is the fish's mouth. Grammatical word-order alternatives are acceptable, but do not leave the ordinal uninflected after this article or omit one of the commands.
% mdparagraph:P178:end

\end{minipage}\par\addvspace{4pt}

\par\begin{minipage}{\linewidth}\raggedright
% mdparagraph:P179:begin
\textbf{20.9.} Thou wilt find a penny; take that, and give (sęllan) for me and for thee.
% mdparagraph:P179:end

% item: EX-VI-20-09
% key: KEY-EX-VI-20-09
% mdparagraph:P180:begin
\textbf{Ðū gemētst pęning; nim ðone, and sęle for mē and for ðē.}
% mdparagraph:P180:end

% mdparagraph:P181:begin
Present \textbf{gemētst} carries future reference. \textbf{Pęning} is masculine; \textbf{ðone} refers back to that penny as the object of \textbf{nim}, and the same object is understood with \textbf{sęle}. \textbf{For mē} and \textbf{for ðē} are the assigned beneficiary phrases. Other grammatical object placement can preserve the meaning. Do not replace singular \textbf{sęle} with a plural command or introduce \textbf{willan} solely because English uses \textbf{wilt}.
% mdparagraph:P181:end

\end{minipage}\par\addvspace{4pt}

\reserve{7\baselineskip}
% mdparagraph:P182:begin
\section*{Cook, Group VI, Exercise 21}
% mdparagraph:P182:end

\par\begin{minipage}{\linewidth}\raggedright
% mdparagraph:P183:begin
\textbf{21.1.} Then said he, She is not dead.
% mdparagraph:P183:end

% item: EX-VI-21-01
% key: KEY-EX-VI-21-01
% mdparagraph:P184:begin
\textbf{Ðā cwæð hē, Hēo nis dēad.}
% mdparagraph:P184:end

% mdparagraph:P185:begin
\textbf{Cwæð} reports a past saying, while \textbf{Hēo nis dēad} is the speaker's direct present statement. The feminine subject takes the selected predicate \textbf{dēad}. \textbf{Then he said, “She is not dead”} gives the intended quotation. Do not automatically change the quoted present to a past or turn it into an indirect clause with a different reporting construction.
% mdparagraph:P185:end

\end{minipage}\par\addvspace{4pt}

\par\begin{minipage}{\linewidth}\raggedright
% mdparagraph:P186:begin
\textbf{21.2.} Then took he her hand, and said, Maiden, arise.
% mdparagraph:P186:end

% item: EX-VI-21-02
% key: KEY-EX-VI-21-02
% mdparagraph:P187:begin
\textbf{Ðā nōm hē hiere hand, and cwæð, Fǣmne, rīs.}
% mdparagraph:P187:end

% mdparagraph:P188:begin
\textbf{Nōm} and \textbf{cwæð} are past narrative verbs. Possessive \textbf{hiere} identifies whose hand he took; the hand is the accusative object. In direct address, \textbf{fǣmne} is the nominative form of the weak feminine noun, followed by singular command \textbf{rīs}. \textbf{Then he took her hand and said, “Maiden, arise”} is the meaning. Do not give the addressed maiden an oblique ending or use plural \textbf{rīsað} for this one person.
% mdparagraph:P188:end

\end{minipage}\par\addvspace{4pt}

\par\begin{minipage}{\linewidth}\raggedright
% mdparagraph:P189:begin
\textbf{21.3.} The stone wall is seven feet broad, and twenty feet high.
% mdparagraph:P189:end

% item: EX-VI-21-03
% key: KEY-EX-VI-21-03
% mdparagraph:P190:begin
\textbf{Se stānweall is seofon fōta brād, and twēntig fōta hēah.}
% mdparagraph:P190:end

% mdparagraph:P191:begin
\textbf{Brād} and \textbf{hēah} agree with masculine \textbf{stānweall}. Both \textbf{seofon fōta} and the supported \textbf{twēntig fōta} express extent, using genitive plural \textbf{fōta}. The selected root-noun plural \textbf{fēt} belongs to other cases and cannot replace \textbf{fōta} here merely because English has \textbf{feet}. Preserve both measurements and their different dimensions; breadth and height are not interchangeable.
% mdparagraph:P191:end

\end{minipage}\par\addvspace{4pt}

\par\begin{minipage}{\linewidth}\raggedright
% mdparagraph:P192:begin
\textbf{21.4.} I took leave of her (acc.), and announced (om. to) her that I never more would visit her.
% mdparagraph:P192:end

% item: EX-VI-21-04
% key: KEY-EX-VI-21-04
% mdparagraph:P193:begin
\textbf{Ic grētte hīe, and cȳðde hiere ðæt ic hīe nǣfre mā genēosian nolde.}
% mdparagraph:P193:end

% mdparagraph:P194:begin
\textbf{Grētte} is past leave-taking with accusative \textbf{hīe}. \textbf{Cȳðde} has the dative person \textbf{hiere} and the following content clause. Within that report, the second \textbf{hīe} is the object of \textbf{genēosian}. The local \textbf{nolde} is the supplied negative preterite subjunctive first singular of \textbf{willan}. With \textbf{nǣfre mā} it expresses one negative intention: never again to visit her. This is prospective from the past announcement, not an assertion that a visit was completed. Other evidenced negative constructions may be valid, but the supplied route requires no independently recalled negative paradigm. Do not turn the announcement recipient into the visit's subject or make negative concord cancel the refusal. [2, 6]
% mdparagraph:P194:end

\end{minipage}\par\addvspace{4pt}

\par\begin{minipage}{\linewidth}\raggedright
% mdparagraph:P195:begin
\textbf{21.5.} The righteous (om. ones) are as bold as (om. a) lion.
% mdparagraph:P195:end

% item: EX-VI-21-05
% key: KEY-EX-VI-21-05
% mdparagraph:P196:begin
\textbf{Ðā rihtwīsan sindon swā bealde swā lēo.}
% mdparagraph:P196:end

% mdparagraph:P197:begin
\textbf{Ðā rihtwīsan} is the articled weak adjective used for righteous people; no noun translating \textbf{ones} is added. Strong predicate \textbf{bealde} agrees with that masculine plural. The first \textbf{swā} marks degree, and the second introduces the comparison standard \textbf{lēo}, supplied as a nominative singular without an indefinite article. The separated \textbf{swā … bealde … swā} construction is the target. Singular \textbf{beald} would fail agreement; adding a noun for \textbf{ones} or an article before \textbf{lēo} would disregard Cook's omission controls. Do not replace the supplied lion form by a guessed declension.
% mdparagraph:P197:end

\end{minipage}\par\addvspace{4pt}

\par\begin{minipage}{\linewidth}\raggedright
% mdparagraph:P198:begin
\textbf{21.6.} The Britons were cowardly.
% mdparagraph:P198:end

% item: EX-VI-21-06
% key: KEY-EX-VI-21-06
% mdparagraph:P199:begin
\textbf{Bryttas wǣron earge.}
% mdparagraph:P199:end

% mdparagraph:P200:begin
Supplied \textbf{Bryttas} is the plural subject; \textbf{wǣron} is preterite plural and \textbf{earge} the agreeing strong masculine plural predicate. Equivalent word order is acceptable. Singular \textbf{wæs}, or a singular adjective used only because English \textbf{cowardly} does not visibly change, would fail agreement.
% mdparagraph:P200:end

\end{minipage}\par\addvspace{4pt}

\par\begin{minipage}{\linewidth}\raggedright
% mdparagraph:P201:begin
\textbf{21.7.} Come hither, fair maiden.
% mdparagraph:P201:end

% item: EX-VI-21-07
% key: KEY-EX-VI-21-07
% mdparagraph:P202:begin
\textbf{Cum hider, fægere fǣmne.}
% mdparagraph:P202:end

% mdparagraph:P203:begin
\textbf{Cum} addresses one person, and \textbf{hider} gives movement hither. Direct address takes weak feminine \textbf{fægere} with nominative \textbf{fǣmne}. The phrase means \textbf{fair maiden}, not a statement that an independently named maiden is fair. \textbf{Come here, fair maiden} is equivalent English. Do not replace the directional adverb with a location adverb without acknowledging the difference, or use a strong adjective ending against the taught address pattern.
% mdparagraph:P203:end

\end{minipage}\par\addvspace{4pt}

\par\begin{minipage}{\linewidth}\raggedright
% mdparagraph:P204:begin
\textbf{21.8.} Happy is that people whose God is the Lord.
% mdparagraph:P204:end

% item: EX-VI-21-08
% key: KEY-EX-VI-21-08
% mdparagraph:P205:begin
\textbf{Ēadig is ðæt folc ðæs God is dryhten.}
% mdparagraph:P205:end

% mdparagraph:P206:begin
\textbf{Ēadig} is the strong neuter nominative predicate agreeing with \textbf{ðæt folc}. Relative \textbf{ðæs} is genitive, whose, because it expresses the possessor of \textbf{God} inside the relative clause. Its case is not copied from the nominative antecedent. Predicate \textbf{dryhten} identifies God as Lord. Equivalent word order preserves the meaning; making \textbf{ðæs} nominative would lose the possessive relation, and a masculine happy-predicate analysis would ignore the noun's grammatical gender.
% mdparagraph:P206:end

\end{minipage}\par\addvspace{4pt}

\par\begin{minipage}{\linewidth}\raggedright
% mdparagraph:P207:begin
\textbf{21.9.} Learn of (æt) me, for I am merciful and lowly (ēaðmōd) in heart, and ye shall find (gemētan) rest (ręst, 51. b) (om. unto) your souls.
% mdparagraph:P207:end

% item: EX-VI-21-09
% key: KEY-EX-VI-21-09
% mdparagraph:P208:begin
\textbf{Leorniað æt mē, for-ðām ic eom mildheort and ēaðmōd on heortan, and gē gemētað ręste ēowrum sāwlum.}
% mdparagraph:P208:end

% mdparagraph:P209:begin
Plural \textbf{leorniað} takes the taught source/\allowbreak{}person phrase \textbf{æt mē}. Causal \textbf{for-ðām} introduces the reason; \textbf{on heortan} expresses the respect in which the speaker is lowly. \textbf{Gemētað} is present with future reference. \textbf{Ręste} is the accusative thing found; \textbf{ēowrum sāwlum} is the dative plural beneficiary, with agreeing possessive. Cook directs omission of a separate translation for \textbf{unto}. Preserve both merciful and lowly, and do not turn the souls into the direct object or add a future auxiliary merely from English \textbf{shall}.
% mdparagraph:P209:end

\end{minipage}\par\addvspace{4pt}

\par\begin{minipage}{\linewidth}\raggedright
% mdparagraph:P210:begin
\textbf{21.10.} Now I send you as sheep among (gemang, 172) wolves.
% mdparagraph:P210:end

% item: EX-VI-21-10
% key: KEY-EX-VI-21-10
% mdparagraph:P211:begin
\textbf{Nū ic sęnde ēow swā scēap gemang wulfas.}
% mdparagraph:P211:end

% mdparagraph:P212:begin
\textbf{Ēow} is the plural accusative person sent. The single \textbf{swā} introduces a comparison phrase, \textbf{as sheep}, with \textbf{scēap} referring to those people; it is not the two-token degree construction of 21.5. The selected motion frame \textbf{gemang + accusative} gives \textbf{wulfas}. Another defensible grammatical order may retain the same relations. Do not treat \textbf{swā} as merely \textbf{so}, require a finite verb after this nominal comparison, or change the wolves into the people being sent.
% mdparagraph:P212:end

\end{minipage}\par\addvspace{4pt}

\par\begin{minipage}{\linewidth}\raggedright
% mdparagraph:P213:begin
\textbf{21.11.} When they pursue (ēhtan) you in this city, flee into (on) (om. an)other.
% mdparagraph:P213:end

% item: EX-VI-21-11
% key: KEY-EX-VI-21-11
% mdparagraph:P214:begin
\textbf{Ðonne hīe ēow ēhtað on ðisse ceastre, flēoð on ōðre.}
% mdparagraph:P214:end

% mdparagraph:P215:begin
\textbf{Ðonne} introduces the temporal circumstance. The model uses the directly evidenced accusative object \textbf{ēow} with \textbf{ēhtað}, as the unit directs. Genitive government is also documented and can be valid language, but does not meet this selected accusative target. \textbf{On ðisse ceastre} is a dative location. Plural command \textbf{flēoð} answers plural \textbf{you}; \textbf{on ōðre} is the locally supplied accusative destination with \textbf{ceastre} understood. Thus \textbf{another} means another city. No separate indefinite article is added. Do not treat both \textbf{on} phrases as static locations or require the learner to reconstruct the supplied ellipse without its help. [6]
% mdparagraph:P215:end

\end{minipage}\par\addvspace{4pt}

\par\begin{minipage}{\linewidth}\raggedright
% mdparagraph:P216:begin
\textbf{21.12.} I take leave of my native land.
% mdparagraph:P216:end

% item: EX-VI-21-12
% key: KEY-EX-VI-21-12
% mdparagraph:P217:begin
\textbf{Ic grēte mīnne ēðel.}
% mdparagraph:P217:end

% mdparagraph:P218:begin
Present first-singular \textbf{grēte} contrasts with past \textbf{grētte} in 21.4. \textbf{Mīnne ēðel} is the masculine accusative native land being saluted in leave-taking. The whole meaning is \textbf{I bid farewell to my native land}; the English preposition \textbf{of} does not impose an Old English genitive here. Other grammatical word orders are acceptable, but a past verb would change the time and a dative object would change the selected frame.
% mdparagraph:P218:end

\end{minipage}\par\addvspace{4pt}

\clearpage
% mdparagraph:P219:begin
\section*{R. A house and a return}
% mdparagraph:P219:end

\reserve{5\baselineskip}
\par\begin{minipage}{\linewidth}\raggedright
% mdparagraph:P220:begin
\textbf{R01. Complete translation.}
% mdparagraph:P220:end

% item: NEW-U12-R01
% key: KEY-NEW-U12-R01
% mdparagraph:P221:begin
The man had built one house. Then his son came, and the son saw the house. The father said, “Now I have built the house.” The son said that he intended to travel. When the son departed, the father sat in the house. Then the son came, and he gave his father food.
% mdparagraph:P221:end

% mdparagraph:P222:begin
\textbf{Interpretation:} The opening \textbf{hæfde … geworht} places the completed building before the son's first arrival. The direct \textbf{ic} belongs to the father; the next reported intention belongs to the son. \textbf{Intended/\allowbreak{}wanted to go} is also a defensible rendering of \textbf{wolde faran}. The following clauses separately narrate departure and return, so fulfillment is established by those events rather than by \textbf{wolde} alone.
% mdparagraph:P222:end

% mdparagraph:P223:begin
\textbf{One house} preserves the assigned cardinal \textbf{ān}; the account has not introduced an indefinite article. Contextual English \textbf{was sitting} can render the father's \textbf{sæt}, although the Old English word is a simple finite preterite. Preserve the son's return and his gift of food to his father. Making the father the giver to himself, or treating the last arrival as the first one, would break the connected reference and sequence. [8]
% mdparagraph:P223:end

\end{minipage}\par\addvspace{4pt}

\par\begin{minipage}{\linewidth}\raggedright
% mdparagraph:P224:begin
\textbf{R02. The order of events.}
% mdparagraph:P224:end

% item: NEW-U12-R02
% key: KEY-NEW-U12-R02
% mdparagraph:P225:begin
\itempara{1}{Building the house.}
% mdparagraph:P225:end

% mdparagraph:P226:begin
\itempara{2}{The son's first arrival.}
% mdparagraph:P226:end

% mdparagraph:P227:begin
\itempara{3}{The father's announcement.}
% mdparagraph:P227:end

% mdparagraph:P228:begin
\itempara{4}{The son's departure.}
% mdparagraph:P228:end

% mdparagraph:P229:begin
\itempara{5}{The son's return with food.}
% mdparagraph:P229:end

% mdparagraph:P230:begin
The phrase \textbf{hæfde … geworht} puts the building before the story's first arrival. Between the announcement and departure the son expresses his intention to travel. A prose description of the same order is equivalent. Do not put building after the announcement merely because the father speaks about it then, or merge the two arrivals.
% mdparagraph:P230:end

\end{minipage}\par\addvspace{4pt}

\par\begin{minipage}{\linewidth}\raggedright
% mdparagraph:P231:begin
\textbf{R03. Speakers, intentions and events.}
% mdparagraph:P231:end

% item: NEW-U12-R03
% key: KEY-NEW-U12-R03
% mdparagraph:P232:begin
\textbf{Ic} in the direct quotation is the father. \textbf{Hē} in \textbf{ðæt hē wolde faran} is the son whose intention is being reported. \textbf{Hē} in the last sentence is the returning son, who gives food to his father.
% mdparagraph:P232:end

% mdparagraph:P233:begin
The journey is initially an intention; \textbf{se sunu gewāt} then explicitly narrates the departure. The subsequent \textbf{cōm se sunu} supplies the return. Thus the narrative shows action beyond the statement of intention, although it gives no detailed account of the intervening travel. Masculine pronoun morphology alone cannot distinguish father from son: the event chain and local subjects identify the referents. Do not infer completed travel solely from \textbf{wolde} or make every \textbf{hē} refer to the same man regardless of context.
% mdparagraph:P233:end

\end{minipage}\par\addvspace{4pt}

\par\begin{minipage}{\linewidth}\raggedright
% mdparagraph:P234:begin
\textbf{R04. The house as passive subject.}
% mdparagraph:P234:end

% item: NEW-U12-R04
% key: KEY-NEW-U12-R04
% mdparagraph:P235:begin
In \textbf{Nū hæbbe ic ðæt hūs geworht}, \textbf{ic} is the nominative subject and builder; \textbf{ðæt hūs} is the accusative object.
% mdparagraph:P235:end

% mdparagraph:P236:begin
\textbf{Ðæt hūs is geworht.} — The house is built.
% mdparagraph:P236:end

% mdparagraph:P237:begin
The new sentence makes \textbf{ðæt hūs} nominative subject and leaves the builder unexpressed. Its completed state is conveyed with present \textbf{is} and the same past-participle base. \textbf{Geworht is ðæt hūs} is a defensible alternative order. The article and noun have identical nominative and accusative spellings, but their grammatical role has changed. Keeping \textbf{ic} as the subject of this agentless passive would not perform the requested transformation.
% mdparagraph:P237:end

\end{minipage}\par\addvspace{4pt}

\clearpage
% mdparagraph:P238:begin
\section*{S. The Picts and Scots}
% mdparagraph:P238:end

\reserve{5\baselineskip}
\reserve{8\baselineskip}
% mdparagraph:P239:begin
\textbf{S01. Complete translation.}
% mdparagraph:P239:end

% item: NEW-U12-S01
% key: KEY-NEW-U12-S01
% mdparagraph:P240:begin
Then two peoples, the Picts from the north and the Scots from the west, began to attack the Britons, take their possessions and plunder; and they afflicted and oppressed them for many years. Then, amid that unrest, the Britons sent messengers to Rome with letters and a tearful petition. They asked the Romans for help and promised them humble obedience and constant submission if the Romans helped them so that they could overcome their enemies. Then the Romans sent them a large army for support; and as soon as the soldiers came onto this island, they fought against their enemies, inflicted great slaughter on them, and drove them out of the Britons' borders and put them to flight. They instructed the Britons to build a fortification for their own protection against their enemies; and so the Romans went home with great victory.
% mdparagraph:P240:end

% mdparagraph:P241:begin
When the former adversaries then realized this, that the Roman army had gone away, they soon came with a fleet into the Britons' territory, struck down and killed all they came across, trampled and destroyed them like ripe crops, and utterly reduced them to misery.
% mdparagraph:P241:end

% mdparagraph:P242:begin
\textbf{Interpretation:} The English names some participants where the Old English uses pronouns: the initial raiders are the Picts and Scots; the petitioners are the Britons; the senders and relieving soldiers are Roman; the final attackers are the former adversaries. Retaining English pronouns is equally acceptable if those relations stay clear. The full sequence includes both the successful relief and the renewed attack after withdrawal.
% mdparagraph:P242:end

% mdparagraph:P243:begin
In opening \textbf{and hīe fela gēara iermdon and hīendon}, \textbf{hīe} can be accusative, the Britons, with the raiders continuing as subject, or nominative, the raiders, with the Britons understood as object. The surface does not require a unique case analysis; either preserves the same oppression. \textbf{Hiera gemǣrum} in the relief describes the Britons' borders, as Cook's note explains.
% mdparagraph:P243:end

% mdparagraph:P244:begin
\textbf{Had gone away} makes the Roman departure earlier than the enemies' realization; \textbf{was gone} can express the same resulting absence. The phrase is intransitive, not passive. \textbf{All they came across} allows the breadth of the neuter headless \textbf{ðæt} clause; \textbf{all whom they encountered} can also make the human victims explicit in context without turning \textbf{ðæt} into an overt masculine relative. \textbf{Ripe crops}, \textbf{ripe corn}, or a suitably understood \textbf{ripe field} can convey the comparison. It describes the Britons' devastation, not an unrelated literal harvest.
% mdparagraph:P244:end

% mdparagraph:P245:begin
Closely related words for aid, oppression, submission, fortification and borders are acceptable. Keep both the instruction and its prospective building; do not add an assertion that the fortification was completed. Keep the final \textbf{mid ealle}, completely or utterly, and the last reduction to misery. The nineteenth-century ethnological speculation in the original apparatus is archival commentary, not part of this translation or a historical conclusion being tested. [7]
% mdparagraph:P245:end

\reserve{5\baselineskip}
\par\begin{minipage}{\linewidth}\raggedright
% mdparagraph:P246:begin
\textbf{S02. Raid, relief and renewed attack.}
% mdparagraph:P246:end

% item: NEW-U12-S02
% key: KEY-NEW-U12-S02
% mdparagraph:P247:begin
\itempara{1}{First raids.}
% mdparagraph:P247:end

% mdparagraph:P248:begin
\itempara{2}{Appeal to Rome.}
% mdparagraph:P248:end

% mdparagraph:P249:begin
\itempara{3}{Roman relief.}
% mdparagraph:P249:end

% mdparagraph:P250:begin
\itempara{4}{Roman withdrawal.}
% mdparagraph:P250:end

% mdparagraph:P251:begin
\itempara{5}{Renewed attack.}
% mdparagraph:P251:end

% mdparagraph:P252:begin
The former adversaries discover that the Roman army has gone away; they then return by fleet. The narrator thus presents the loss of Roman protection as the opportunity for the renewed attack. Equivalent wording or a labeled timeline is acceptable. The withdrawal precedes its discovery, and the first relief genuinely occurs within the narrative; neither should be omitted or placed after the renewed attack.
% mdparagraph:P252:end

\end{minipage}\par\addvspace{4pt}

\par\begin{minipage}{\linewidth}\raggedright
% mdparagraph:P253:begin
\textbf{S03. Three occurrences of ðā.}
% mdparagraph:P253:end

% item: NEW-U12-S03
% key: KEY-NEW-U12-S03
% mdparagraph:P254:begin
\textbf{Ðā ðæt ðā ongēaton ðā ǣrran gewinnan} can be unpacked as \textbf{When the former adversaries then realized this}.
% mdparagraph:P254:end

% mdtable:T11:begin
\par\begingroup
\fontsize{10.5}{13.2}\selectfont
\setlength{\LTleft}{0pt}\setlength{\LTright}{\fill}
\noindent
\begin{tabular}{@{}>{\raggedright\arraybackslash}p{25.2mm}>{\raggedright\arraybackslash}p{95.99mm}@{}}
\toprule
\textbf{Form in order} & \textbf{Function} \\
\midrule
First \textbf{Ðā} & Temporal conjunction, when. \\
Second \textbf{ðā} & Temporal adverb, then, within that clause. \\
Third \textbf{ðā} & Nominative plural article with \textbf{ǣrran gewinnan}, the former adversaries. \\
\bottomrule
\end{tabular}
\par\endgroup

% mdtable:T11:end
% mdparagraph:P255:begin
The first \textbf{ðæt} is an anticipatory accusative object of \textbf{ongēaton}, this or that. Its content is expanded by the following \textbf{ðæt se Rōmanisca hęre …} clause. The later main clause begins \textbf{ðā cōmon hīe}, then they came. English may omit an unnecessary \textbf{then} while preserving the structure. Do not classify all three \textbf{ðā} forms as articles, or confuse the anticipatory object with the following content-clause conjunction. Cook's original note supplies the when–then–the distinction. [7]
% mdparagraph:P255:end

\end{minipage}\par\addvspace{4pt}

\par\begin{minipage}{\linewidth}\raggedright
% mdparagraph:P256:begin
\textbf{S04. Changing subjects.}
% mdparagraph:P256:end

% item: NEW-U12-S04
% key: KEY-NEW-U12-S04
% mdparagraph:P257:begin
The \textbf{Britons} send messengers to Rome. The \textbf{Romans} send the army to the Britons. The \textbf{Roman soldiers} arrive and fight the enemies.
% mdparagraph:P257:end

% mdparagraph:P258:begin
The first act is the appeal of the afflicted people. The second is Rome's response. In the arrival and fighting clauses the Roman soldiers supply the contextual subject; Cook's note calls them the legionaries, apparently. His next note identifies the borders from which the enemies are driven as the Britons' borders. Repeated \textbf{hīe} does not by itself name a people or guarantee that one subject persists through every clause. \textbf{The Romans} or \textbf{the Roman authorities} is sufficient for the sending: no particular ruler has to be invented. [7]
% mdparagraph:P258:end

\end{minipage}\par\addvspace{4pt}

\par\begin{minipage}{\linewidth}\raggedright
% mdparagraph:P259:begin
\textbf{S05. Wæs onweg gewiten.}
% mdparagraph:P259:end

% item: NEW-U12-S05
% key: KEY-NEW-U12-S05
% mdparagraph:P260:begin
\textbf{Wæs} is third-singular preterite \textbf{wesan}, agreeing with \textbf{se Rōmanisca hęre}. \textbf{Onweg} means away. \textbf{Gewiten} is the past participle of intransitive \textbf{gewītan}.
% mdparagraph:P260:end

% mdparagraph:P261:begin
The army's departure precedes the former enemies' realization that it has left, and their renewed attack follows that realization. \textbf{Had gone away} or \textbf{was gone} preserves the relation. The army itself has departed; it is not an object on which an unnamed agent has performed a transitive action. Do not analyze the phrase as an agentless passive merely because it contains \textbf{wesan} and a past participle, or place the departure after the renewed invasion.
% mdparagraph:P261:end

\end{minipage}\par\addvspace{4pt}

\par\begin{minipage}{\linewidth}\raggedright
% mdparagraph:P262:begin
\textbf{S06. Active Old English and an English passive.}
% mdparagraph:P262:end

% item: NEW-U12-S06
% key: KEY-NEW-U12-S06
% mdparagraph:P263:begin
In \textbf{Ðā onsęndon hīe him micelne hęre}, \textbf{hīe}, the Romans, is nominative subject; \textbf{him}, the Britons, is the dative recipient; and \textbf{micelne hęre} is the accusative army sent.
% mdparagraph:P263:end

% mdparagraph:P264:begin
\textbf{Then a large army was sent to the Britons by the Romans for their support.}
% mdparagraph:P264:end

% mdparagraph:P265:begin
This English passive makes the army the grammatical subject and retains the recipient and the contextually identified agent. \textbf{A large army was then sent to them to help them} is also apt when the Roman sender is clear from context. Cook compares an active Old English clause with a passive Latin source; \textbf{onsęndon} itself remains active. No Latin analysis or new Old English agent phrase is required. Making the Britons the army being sent would confuse recipient and object. [7]
% mdparagraph:P265:end

\end{minipage}\par\addvspace{4pt}

\par\begin{minipage}{\linewidth}\raggedright
% mdparagraph:P266:begin
\textbf{S07. A prospective instruction.}
% mdparagraph:P266:end

% item: NEW-U12-S07
% key: KEY-NEW-U12-S07
% mdparagraph:P267:begin
The Romans recommend that the Britons build a fortification for their own protection against their enemies. The building is \textbf{prospective} from the instruction, and this clause alone does not state that the work was completed.
% mdparagraph:P267:end

% mdparagraph:P268:begin
\textbf{Lǣrdon} reports the instruction in the past. Dependent \textbf{worhten} is preterite subjunctive in that reporting setting; its morphology does not make the recommended building earlier than the advice. \textbf{Instructed/\allowbreak{}advised them to build} conveys the same relation. Translating it as \textbf{they had already built} merely because the dependent form is preterite would reverse the task's temporal meaning. Cook's note about construction in an omitted Latin passage does not turn this Old English advice clause into a narrated completion. [7]
% mdparagraph:P268:end

\end{minipage}\par\addvspace{4pt}

\par\begin{minipage}{\linewidth}\raggedright
% mdparagraph:P269:begin
\textbf{S08. The crop comparison.}
% mdparagraph:P269:end

% item: NEW-U12-S08
% key: KEY-NEW-U12-S08
% mdparagraph:P270:begin
The victims are the Britons encountered by the returning attackers. \textbf{Rīpe ierðe}, ripe crops, is the supplied collective phrase used in comparing their treatment to a crop trampled down and destroyed. The image intensifies the narrative's account of devastation and vulnerability after the Roman departure.
% mdparagraph:P270:end

% mdparagraph:P271:begin
\textbf{Ripe corn} or an appropriately understood \textbf{ripe field} can also convey the image. A broader rendering of the preceding \textbf{all they came across} remains possible; the grammatical neuter \textbf{ðæt} must not be converted into a uniquely personal relative merely to match an English model. The comparison does not begin a separate literal farming episode or give the noun for crops a newly taught ordinary meaning of people. Making the Romans the victims would disregard their departure and the changed actors. [7]
% mdparagraph:P271:end

\end{minipage}\par\addvspace{4pt}

\clearpage
% mdparagraph:P272:begin
\section*{Source notes}
% mdparagraph:P272:end

\begingroup\fontsize{10}{12.8}\selectfont
% mdparagraph:P273:begin
[1] Cook, \emph{Exercises in Old English} (1895), Group VI, Exercises 19–21, printed pp. 9–11; one-based PDF pp. 17–19. All 38 original prompts, controls and section references are retained. In 20.6 the checked hint is \textbf{gīet}; in 21.11 the printed adjacency \textbf{(om. an)other.} is retained. All answers and explanatory prose are new teaching material.
% mdparagraph:P273:end

% mdparagraph:P274:begin
[2] Cook, \emph{A First Book in Old English}, second edition (1895), §§186–188, printed pp. 101–102; Fulk, \emph{An Introductory Grammar of Old English} (2020), §§45–46 and 48–52, printed pp. 28–29 and 31–33. These distinguish the two inflected tenses, contextual temporal relations, reported presentation, compound constructions and the limits of an automatic modern progressive analysis. The present and past reporting examples retain the task's stated simultaneous, anterior or intended-later relation.
% mdparagraph:P274:end

% mdparagraph:P275:begin
[3] Cook, \emph{First Book}, §188, printed p. 102; Bosworth–Toller, \textbf{habban}, main dictionary, entry beginning p. 496b, especially IV–V on p. 497a, and supplement, auxiliary subsection (1)(a–b) on p. 493a. The directly inspected entries distinguish uninflected and object-agreeing participles. The selected active compound pattern does not assign every causative or resultative construction appearing elsewhere in those entries.
% mdparagraph:P275:end

% mdparagraph:P276:begin
[4] Cook, \emph{First Book}, the checked noun patterns on printed pp. 28–38, including §44.2 for the \textbf{hęre} alternatives and §46 for the root nouns, with the relevant lexical entries. Complete selected forms are printed in F01–F02. Documented alternatives apply to their identified cells and do not replace the required paradigm with an unrestricted choice of stems.
% mdparagraph:P276:end

% mdparagraph:P277:begin
[5] Cook, \emph{First Book}, the strong and weak verb systems on printed pp. 54–78, with the checked lexical entries; especially §93 for the strong second-singular preterite stem, the contracted strong forms on pp. 58–61, and the weak systems on pp. 68–74. Wright, \emph{Old English Grammar} (1914), §476, printed p. 248 /\allowbreak{} one-based PDF p. 268, supports the selected dental simplification in \textbf{cȳst/\allowbreak{}cȳð}; \textbf{cȳðst} follows the documented analogical retained-dental pattern and remains a valid alternative. This is a rule-based formation, not a claim that these checked pages attest the exact lexical form. These references establish morphology within the defined systems, not every possible use of a displayed cell.
% mdparagraph:P277:end

% mdparagraph:P278:begin
[6] Cook, \emph{First Book}, §139, printed p. 79, supplies negative \textbf{nolde} in the indicative and subjunctive systems; 21.4 uses the locally supplied first-singular preterite subjunctive. Bosworth–Toller, \textbf{ēhtan}, main p. 244, records both genitive and accusative government and the directly relevant accusative example; \textbf{gemang}, main p. 414, records the selected accusative sending-among construction. The separate \textbf{swā} comparisons are supported by the main entry on p. 940; numerical predicative \textbf{mā} is supported by the main entry on p. 653. These are specific constructions, not general case rules inferred from English prepositions.
% mdparagraph:P278:end

% mdparagraph:P279:begin
[7] Cook, \emph{First Book}, Reader VI, printed p. 145, editorial prose lines 1–7, and p. 146, editorial prose lines 1–11 through \textbf{and hīe mid ealle foriermdon.}; one-based PDF pp. 161–162. The assignment ends before \textbf{And hīe ðā eft sęndon}. Its 21 attached notes are p. 145 notes 1–6 and p. 146 notes 1–15. On p. 146, notes 3–5 explain active/\allowbreak{}passive wording, legionaries and the Britons' borders; notes 10–11 address the omitted Latin construction passage and the three \textbf{ðā} forms; note 13 reads \textbf{Lat. cædunt} at both anchors; notes 14–15 support the crop phrase and \textbf{mid ealle}. Original apparatus and new lexical help remain distinct. The full translation was independently construed from the Old English before comparison with the other editorial witnesses; equivalent interpretations are not restricted to its English wording.
% mdparagraph:P279:end

% mdparagraph:P280:begin
[8] \emph{A house and a return}, all original F/\allowbreak{}A/\allowbreak{}C/\allowbreak{}R/\allowbreak{}S tasks and the twelve teaching examples are new composition for this course. The account uses the contracted productive vocabulary and constructions; it is not represented as an ancient quotation or as an adaptation of the Picts and Scots passage.
% mdparagraph:P280:end

\endgroup
\end{document}
